%%%%%%%%%%%%%%%%%%%%%%%%%%%%%%%%%%%%%%%%%%%%%%%%%%%%%%%%%%%%%%%%%%%%%%
% Lớp 12 | Phương trình mặt phẳng, đường thẳng, mặt cầu | Phương trình mặt phẳng
% Dạng 2H5?1-6: Góc giữa hai mặt phẳng
% Số câu: 33  (sinh tự động bởi tach_cau.py; từ nay thầy sửa trực tiếp file này)
%%%%%%%%%%%%%%%%%%%%%%%%%%%%%%%%%%%%%%%%%%%%%%%%%%%%%%%%%%%%%%%%%%%%%%

%%%%%%%%%%%%%%%% TRẮC NGHIỆM 4 PHƯƠNG ÁN %%%%%%%%%%%%%%%%

% nguồn: Phan1_Lop_12_Hinh_hoc_Chuong_5
\begin{ex}%[2H5N1-6]%[ma:ff969214]
	Trong không gian $Oxyz$,  cho hai mặt phẳng $(\alpha): x + y - 2z + 1 = 0$ và $(\beta): x + 2y + 1 = 0$. Tính côsin góc giữa hai mặt phẳng $(\alpha)$ và $(\beta)$.
	 \choice
		 {$\dfrac{\sqrt{10}}{10}$}
		{\True  $\dfrac{\sqrt{30}}{10}$}
		 {$\dfrac{1}{6}$}
		{$\dfrac{1}{3}$} 
	\loigiai{ 
		Ta có $\overrightarrow{n}_\alpha = (1,1,-2)$, $\overrightarrow{n}_\beta = (1,2,0)$ là các vectơ pháp tuyến  của hai mặt phẳng $(\alpha)$ và $(\beta)$.\\
		$\cos((\alpha);(\beta)) = \dfrac{|\overrightarrow{n}_\alpha \cdot \overrightarrow{n}_\beta|}{|\overrightarrow{n}_\alpha|\cdot|\overrightarrow{n}_\beta|} = \dfrac{|1\cdot1 + 1\cdot2 + (-2)\cdot0|}{\sqrt{1^2 + 1^2 + (-2)^2}\cdot\sqrt{1^2  + 2^2}} = \dfrac{3}{\sqrt{6}\cdot\sqrt{5}} = \dfrac{3}{\sqrt{30}} = \dfrac{\sqrt{30}}{10}$.
	 }
\end{ex}

% nguồn: Phan1_Lop_12_Hinh_hoc_Chuong_5
\begin{ex}%[2H5N1-6]%[ma:e8f694ae]
	Trong không gian với hệ trục tọa độ $Oxyz$, góc giữa hai mặt phẳng $(Oxy)$ và  $(Oxz)$ bằng
	 \choice
		 {$45^\circ$}
		 {$30^\circ$}
		 {\True $90^\circ$}
		{$60^\circ$} 
	\loigiai{ 
		Góc giữa hai mặt  phẳng $Oxy$ và $Oxz$ bằng $90^\circ$.
	 }
\end{ex}

% nguồn: Phan1_Lop_12_Hinh_hoc_Chuong_5
\begin{ex}%[2H5H1-6]%[ma:fe9a4135]
	Trong không gian $Oxyz$, góc giữa hai mặt  phẳng $(P) \colon x+2y+z-1=0$ và $(Q) \colon -x+y+2z+2=0$ bằng
	\choice 
		 {$30^\circ$}
		{$45^\circ$} 
		{\True $60^\circ$} 
		 {$90^\circ$}
	 \loigiai{
		Mặt phẳng $(P)$ có một vectơ pháp  tuyến là $\vec{n}_{1} = (1;2;1)$.\\
		 Mặt phẳng $(Q)$ có một vectơ pháp tuyến là $\vec{n}_{2} =(-1;1;2)$.\\
		Gọi $\alpha$ ($0^\circ \leq \alpha \leq 90^\circ$) là góc giữa hai mặt phẳng  $(P)$ và $(Q)$.\\Ta có\\
		$\cos(\alpha) = \dfrac{|\vec{n}_{1}\cdot \vec{n}_{2}|}{|\vec{n}_{1}|\cdot|\vec{n}_{2}|} = \dfrac{|1\cdot (-1)+2\cdot 1+1 \cdot 2|}{\sqrt{1^2 + 2^2+1^2}\cdot \sqrt{(-1)^2  +1^2 +2^2}} = \dfrac{3}{\sqrt{6}\cdot \sqrt{6}} = \dfrac{3}{6} = \dfrac{1}{2} \Rightarrow \alpha = 60^\circ$.
	} 
\end{ex}

% nguồn: Phan1_Lop_12_Hinh_hoc_Chuong_5
\begin{ex}%[2H5H1-6]%[ma:6ec1ae87]
	Trong không gian  $Oxyz$, cho hai mặt phẳng $(P)$ và $(Q)$ lần lượt có hai vectơ pháp tuyến là $\overrightarrow{n}_P$ và $\overrightarrow{n}_Q$. Biết góc giữa hai vectơ $\overrightarrow{n}_P$ và $\overrightarrow{n}_Q$ bằng $120^\circ$. Góc giữa hai mặt phẳng $(P)$ và $(Q)$ bằng
	 \choice
		{\True  $60^\circ$}
		 {$120^\circ$}
		{$90^\circ$} 
		{$45^\circ$} 
	\loigiai{ 
		Gọi $\varphi$ là góc giữa hai mặt phẳng $(P)$ và  $(Q)$.\\
		Do  $\left(\overrightarrow{n}_P, \overrightarrow{n}_Q\right)=120^\circ > 90^\circ \Rightarrow \varphi=180^\circ-120^\circ=60^\circ$.\\
		Vậy góc giữa hai mặt phẳng $(P)$ và $(Q)$ bằng $60^\circ$. 
	 }
\end{ex}

% nguồn: Phan1_Lop_12_Hinh_hoc_Chuong_5
\begin{ex}%[2H5H1-6]%[ma:3ad44bcd]
	Trong không gian $Oxyz$, cho hai mặt phẳng $(P)$ và $(Q)$ lần lượt  có hai vectơ pháp tuyến là $\overrightarrow{n}_P$ và $\overrightarrow{n}_Q$. Biết góc giữa hai vectơ $\overrightarrow{n}_P$ và $\overrightarrow{n}_Q$ bằng $120^\circ$. Góc giữa hai mặt phẳng $(P)$ và $(Q)$ bằng
	\choice 
		 {$45^\circ$}
		 {$120^\circ$}
		 {$30^\circ$}
		 {\True $60^\circ$}
	\loigiai{ 
		Ta có $\left(\widehat{\overrightarrow{n}_P, \overrightarrow{n}_Q}\right)=120^\circ  \Rightarrow(\widehat{(P),(Q)})=60^\circ$.
	} 
\end{ex}

% nguồn: Phan1_Lop_12_Hinh_hoc_Chuong_5
\begin{ex}%[2H5H1-6]%[ma:aad1d428]
	Trong không gian với hệ trục tọa độ $Oxyz$, cho hai mặt phẳng $(P)\colon x-2y+2z-1=0$  và $(Q)\colon 2x+2y-z-3=0$. Gọi $\alpha$ là góc giữa hai mặt phẳng $(P)$ và $(Q)$. Giá trị $\cos \alpha$ bằng
	\choice 
		 {$\dfrac{2}{3}$}
		{$-\dfrac{4}{9}$} 
		{$-\dfrac{2}{3}$} 
		{\True  $\dfrac{4}{9}$}
	\loigiai{ 
		Ta có  $\overrightarrow{n}_{(P)} =(1;-2;2)$, $\overrightarrow{n}_{(Q)}=(2;2;-1)$.\\
		Khi đó $\cos  \alpha=\dfrac{\left|\overrightarrow{n}_{(P)} \cdot \overrightarrow{n}_{(Q)}\right|}{\left|\overrightarrow{n}_{(P)}\right|\cdot \left|\overrightarrow{n}_{(Q)}\right|}=\dfrac{|1\cdot2+(-2)\cdot 2+2\cdot (-1)|}{\sqrt{1^2 +(-2)^2+2^2} \cdot \sqrt{2^2+2^2+(-1)^2}}=\dfrac{4}{9}$.
	} 
\end{ex}

% nguồn: Phan1_Lop_12_Hinh_hoc_Chuong_5
\begin{ex}%[2H5H1-6]%[ma:773b6fbc]
	Trong không gian $Oxyz$, côsin của góc giữa hai mặt phẳng $\left(P\right)\colon x+2y-2z+1=0$ và $\left(Q\right)\colon x+y+z-1=0$ bằng 
	\choice 
		{$\dfrac{\sqrt{3}}{3}$} 
		{\True  $\dfrac{\sqrt{3}}{9}$}
		{$-\dfrac{\sqrt{3}}{9}$} 
		 {$-\dfrac{\sqrt{3}}{3}$}
	 \loigiai{
		Mặt phẳng $\left(P\right)\colon x+2y-2z+1=0$ có véc-tơ pháp tuyến là $\vec{n}_{\left(P\right)}=\left(1;2;-2\right)$.\\ 
		Mặt  phẳng $\left(Q\right)\colon x+y+z-1=0$ có véc-tơ pháp tuyến là $\vec{n}_{\left(Q\right)}=\left(1;1;1\right)$.\\
		Ta có $\cos\left(\left(P\right),\left(Q\right)\right)=\left|\cos\left(\vec{n}_{\left( P\right)},\vec{n}_{\left(Q\right)  }\right)\right|=\dfrac{\left|1\cdot1+2\cdot1-2\cdot1\right|}{\sqrt{1^{2}+2^{2}+\left(-2\right)^{2}}\cdot\sqrt{1^{2}+1^{2}+1^{2}}}=\dfrac{\sqrt{3}}{9}$.
	 }
\end{ex}

% nguồn: Phan1_Lop_12_Hinh_hoc_Chuong_5
\begin{ex}%[2H5H1-6]%[ma:f79527d4]
	Trong không gian với hệ trục tọa độ $Oxyz$, cho hai  mặt phẳng $(P)\colon x-2y+2z-1=0$ và $(Q)\colon 2x+2y-z-3=0$. Gọi $\alpha$ là góc giữa hai mặt phẳng $(P)$ và $(Q)$. Tính $\cos \alpha$.
	 \choice
		 {$-\dfrac{4}{9}$}
		{\True $\dfrac{4}{9}$} 
		{$\dfrac{2}{3}$} 
		 {$-\dfrac{2}{3}$}
	\loigiai{ 
		Ta có $\overrightarrow{n}_P=(1;-2; 2)$ và vectơ  pháp tuyến $\overrightarrow{n}_Q=(2; 2;-1)$.\\
		Suy ra $\cos  \alpha=\dfrac{\left|\overrightarrow{n}_P \cdot \overrightarrow{n}_Q\right|}{\left|\overrightarrow{n}_P\right|\cdot\left|\overrightarrow{n}_Q\right|}=\dfrac{|1\cdot 2+(-2) \cdot 2+2\cdot(-1)|}{\sqrt{1^2+(-2)^2+2^2} \cdot \sqrt{2^2+2^2+(-1)^2}}=\dfrac{4}{9}$.
	 }
\end{ex}

% nguồn: Phan1_Lop_12_Hinh_hoc_Chuong_5
\begin{ex}%[2H5H1-6]%[ma:5c6ba605]
	Trong không gian $Oxyz$, côsin của góc giữa hai mặt phẳng $(P)\colon  x+2y-2z+1=0$ và $(Q)\colon x+y+z-1=0$ bằng
	 \choice
		 {$\dfrac{\sqrt{3}}{3}$}
		{\True  $\dfrac{\sqrt{3}}{9}$}
		{$-\dfrac{\sqrt{3}}{9}$} 
		{$-\dfrac{\sqrt{3}}{3}$} 
	\loigiai{ 
		Mặt phẳng $(P)\colon x+2y-2z+1=0$ có véc-tơ pháp tuyến là  $\overrightarrow{n}_{(P)}=\left(1;2;-2\right)$.\\
		Mặt phẳng  $(Q)\colon x+y+z-1=0$ có véc-tơ pháp tuyến là $\overrightarrow{n}_{(Q)}=\left(1;1;1\right)$.\\
		 Ta có $\cos\left((P),(Q)\right)=\left|\cos\left(\overrightarrow{n}_{\left( P\right)},\overrightarrow{n}_{(Q) }\right)\right|=\dfrac{\left|1\cdot1+2\cdot1-2\cdot1\right|}{\sqrt{1^2+2^2+\left(-2\right)^2}\cdot\sqrt{1^2+1^2+1^2}}=\dfrac{\sqrt{3}}{9}$.
	} 
\end{ex}

% nguồn: Phan1_Lop_12_Hinh_hoc_Chuong_5
\begin{ex}%[2H5H1-6]%[ma:a12c131d]
	Trong không gian $Oxyz$, góc giữa hai mặt phẳng $(Oxy)$ và $(Oy z)$  bằng
	 \choice
		 {$30^\circ$}
		 {$45^\circ$}
		 {$60^\circ$}
		{\True  $90^\circ$}
	 \loigiai{
		Ta có vectơ pháp tuyến của $(Oxy)$ và  $(Oy z)$ lần lượt là $\overrightarrow{k}$ và $\overrightarrow{i}$.\\
		Vì $\overrightarrow{k} \perp \overrightarrow{i}$  nên $(\widehat{(Oxy);(Oy z)})=90^\circ$.
	 }
\end{ex}

% nguồn: Phan1_Lop_12_Hinh_hoc_Chuong_5
\begin{ex}%[2H5H1-6]%[ma:7b27cdfe]
	Trong không gian với hệ tọa độ $Oxyz$, góc giữa  hai mặt phẳng $(Ox z)$ và $(Oy z)$ bằng
	 \choice
		{\True  $90^\circ$}
		{$60^\circ$} 
		{$30^\circ$} 
		 {$45^\circ$}
	 \loigiai{
		Dễ thấy, do $Ox$; $Oy$; $Oz$ đôi một vuông góc nên $(Oxz) \perp(Oyz)$.\\ 
		Vậy góc giữa hai  mặt phẳng $(Oxz)$ và $(Oyz)$ bằng $90^\circ$.
	} 
\end{ex}

% nguồn: Phan1_Lop_12_Hinh_hoc_Chuong_5
\begin{ex}%[2H5H1-6]%[ma:7681fd49]
	Cho hai mặt phẳng $(P)\colon 2 x-y-z-3=0$ và $(Q)\colon x-z-2=0$. Góc giữa hai mặt phẳng $(P)$ và $(Q)$ bằng 
	 \choice
		{$30^{\circ}$} 
		{$45^{\circ}$} 
		 {\True $60^{\circ}$}
		{$90^{\circ}$} 
	 \loigiai{
		$(P)$ có vectơ pháp  tuyến là $\overrightarrow{n}_1=(2; -1; -1)$.\\
		$(Q)$ có vectơ pháp tuyến là $\overrightarrow{n}_2=(1; 0; -1)$.\\ 
		Gọi $\varphi$ là góc  giữa hai mặt phẳng $(P)$ và $(Q)$ ta có $$\cos\varphi=\dfrac{\left| 2\cdot 1-1\cdot 0+ 1\cdot 1\right| }{\sqrt{2^2+(-1)^2+(-1)^2}\cdot\sqrt{1^2+0^2+(-1)^2}}=\dfrac{\sqrt{3}}{2}\Rightarrow \varphi=60^{\circ}$$.
	} 
\end{ex}

% nguồn: Phan1_Lop_12_Hinh_hoc_Chuong_5
\begin{ex}%[2H5H1-6]%[ma:3d446025]
	Trong không gian với hệ tọa độ $Oxyz$, cho mặt phẳng $(P)\colon x+y+4z-1=0$ và $(Q)\colon x-z=0$. Góc giữa  hai mặt phẳng $(P)$ và $(Q)$ bằng
	 \choice
		 {$0^\circ$}
		 {$30^\circ$}
		{$45^\circ$} 
		{\True  $60^\circ$}
	\loigiai{ 
		Véctơ pháp tuyến của mặt phẳng $(P)\colon x+y+4z-1=0$ là $\overrightarrow{n}_P  = (1; 1; 4)$.\\
		 Véctơ pháp tuyến của mặt phẳng $(Q)\colon x-z=0$ (hay $x+0y-z=0$) là $\overrightarrow{n}_Q = (1; 0; -1)$.\\
		Gọi $\alpha$ là góc giữa hai mặt phẳng $(P)$ và $(Q)$. Khi đó,  $\cos \alpha$ được tính bằng công thức:
		$$ \cos \alpha  = \dfrac{|\overrightarrow{n}_P \cdot \overrightarrow{n}_Q|}{|\overrightarrow{n}_P| \cdot |\overrightarrow{n}_Q|}=\dfrac{|1\cdot 1+1\cdot 0+4\cdot (-1)|}{\sqrt{1^2+1^2+4^2}\cdot \sqrt{1^2+0^2+(-1)^2}}=\dfrac{3}{3\sqrt{2}\cdot \sqrt{2}}=\dfrac{1}{2}.$$\\
		Vì $0^\circ \le \alpha \le 90^\circ$, nên $\alpha = \arccos\left(\dfrac{1}{2}\right) = 60^\circ$.\\ 
	} 
\end{ex}

% nguồn: Phan1_Lop_12_Hinh_hoc_Chuong_5
\begin{ex}%[2H5H1-6]%[ma:8bf74174]
	Trong không gian với hệ trục $Oxyz$, gọi  $\alpha$ là góc giữa hai mặt phẳng $(P)\colon x+2y-z+2=0$ và $(Q)\colon 2x-y-z+4=0$. Tính $\cos \alpha$.
	 \choice
		{$\cos  \alpha=\dfrac{2}{3}$}
		{$\cos  \alpha=\dfrac{3}{4}$}
		 {\True $\cos \alpha=\dfrac{1}{6}$}
		{$\cos  \alpha=\dfrac{1}{3}$}
	\loigiai{ 
		Mặt phẳng $(P)$ có vectơ pháp tuyến $\overrightarrow{n}_1=(1;  2;-1)$.\\
		Mặt phẳng $(Q)$ có vectơ pháp tuyến  $\overrightarrow{n}_2=(2;-1;-1)$.\\
		Vậy $\cos \alpha=\dfrac{\left|\overrightarrow{n}_1 \cdot \overrightarrow{n}_2\right|}{\left|\overrightarrow{n}_1\right|\cdot\left|\overrightarrow{n}_2\right|}=\dfrac{|2-2+1|}{\sqrt{6} \cdot \sqrt{6}}=\dfrac{1}{6}$. 
	 }
\end{ex}

% nguồn: Phan1_Lop_12_Hinh_hoc_Chuong_5
\begin{ex}%[2H5H1-6]%[ma:e0a249bf]
	Trong không gian với hệ trục  tọa độ $Oxyz$, góc giữa hai mặt phẳng $(P)\colon x+y-z-11=0$ và $(Q)\colon 2x+2y-2z+7=0$ bằng
	 \choice
		 {\True $0^\circ$}
		 {$45^\circ$}
		{$180^\circ$} 
		 {$90^\circ$}
	\loigiai{ 
		Ta có  $\overrightarrow{n}_{(P)}=(1; 1;-1), \overrightarrow{n}_{(Q)}=(2; 2;-2) \Rightarrow \overrightarrow{n}_{(Q)}=2\overrightarrow{n}_{(P)}$.\\
		Do  đó góc giữa hai mặt phẳng $(P)$ và $(Q)$ là $0^\circ$.
	 }
\end{ex}

% nguồn: Phan1_Lop_12_Hinh_hoc_Chuong_5
\begin{ex}%[2H5H1-6]%[ma:96d72001]
	Trong không gian $Oxyz$, góc giữa hai mặt phẳng $(Ox z)$  và $(P)\colon x-y+1=0$ bằng
	\choice 
		{$60^\circ$} 
		{$135^\circ$} 
		{\True $45^\circ$} 
		 {$90^\circ$}
	\loigiai{ 
		Ta có vectơ  pháp tuyến của mặt phẳng $(Ox z)$ là $\overrightarrow{j}=(0; 1; 0)$.\\
		vectơ pháp tuyến của mặt phẳng $(P)\colon x-y+1=0$ là $\overrightarrow{n}=(1;-1;  0)$.\\
		 Gọi $\alpha$ là góc giữa hai mặt phẳng $(Ox z)$ và $(P)$, ta có
		\[\cos \alpha=\dfrac{|\overrightarrow{n} \cdot \overrightarrow{j}|}{|\overrightarrow{n}|\cdot|\overrightarrow{j}|}=\dfrac{|1\cdot 0+1\cdot(-1)+0\cdot 0|}{\sqrt{1^2+(-1)^2+0^2} \cdot  \sqrt{0^2+1^2+0^2}}=\dfrac{1}{\sqrt{2}}. \]
		Vậy  $\alpha=45^\circ$.
	} 
\end{ex}

% nguồn: Phan1_Lop_12_Hinh_hoc_Chuong_5
\begin{ex}%[2H5V1-6]%[ma:13610a37]
	Trong không gian cho hai điểm $E(3;0;0)$, $F(0;0;1)$. Cặp mặt phẳng nào dưới đây đi qua $E$, $F$ đồng thời  tạo với trục $Oz$ một góc $30^{\circ}$?
	\choice 
		{\True $(Q_{1})\colon x-\sqrt{26}y+3z-3=0$ và $(Q_{2})\colon  x+\sqrt{26}y+3z-3=0$}
		{$(Q_{1})\colon x-2\sqrt{6}y+3z-3=0$ và $(Q_{2})\colon x+\sqrt{26}y+3z-3=0$} 
		{$(Q_{1})\colon x-\sqrt{26}y+3z-3=0$ và $(Q_{2})\colon  x+2\sqrt{6}y+3z-3=0$}
		{$(Q_{1})\colon x-4\sqrt{3}y+3z-3=0$  và $(Q_{2})\colon x+3\sqrt{7}y+3z-3=0$}
	 \loigiai{
		Mặt phẳng $(Q)$ đi qua $E$, $F$ nên có dạng $\dfrac{x}{3}+\dfrac{y}{b}+\dfrac{z}{1}=1 \Leftrightarrow bx+3y+3bz-3b=0$, trong đó  $b\neq 0$.\\
		 Khi đó $\overrightarrow{n}_Q=(b;3;3b)$.\\
		Trục $Oz$ có vectơ chỉ phương $\overrightarrow{k}=(0;0;1)$.\\ 
		Góc giữa $(Q)$ và $Oz$ bằng $30^{\circ}$ nên 
		 $$
		 \sin30^\circ=\left|\cos\left(\overrightarrow{n}_{Q},\overrightarrow{k}\right)\right|=\dfrac{|3b|}{\sqrt{b^{2}+9+9b^{2}}}=\dfrac{1}{2}.
		$$ 
		Suy  ra $b=\pm \dfrac{3}{\sqrt{26}}$.\\
		Thay vào ta được  hai mặt phẳng $(Q_{1})\colon x-\sqrt{26}y+3z-3=0$, $(Q_{2})\colon x+\sqrt{26}y+3z-3=0$.
	 }
\end{ex}

% nguồn: Phan1_Lop_12_Hinh_hoc_Chuong_5
\begin{ex}%[2H5V1-6]%[ma:e551ac20]
	Cho hình  lăng trụ tam giác đều $ABC.A'B'C'$ có tất cả các cạnh bằng nhau. Gọi $\alpha$ là góc tạo bởi hai mặt phẳng $(C'AB)$ và $(BCC'B')$. Giá trị $\tan\alpha$ bằng
	 \choice
		 {\True $\sqrt{6}$}
		{$\sqrt{2}$} 
		 {$\dfrac{\sqrt{6}}{2}$}
		 {$\dfrac{2\sqrt{3}}{3}$}
	 \loigiai{
		\begin{center}
			\begin{tikzpicture}[line join = round,line cap=round,>=stealth,font=\footnotesize,scale=1]
					 \path
						(0,2) coordinate  (A)
						(4,2) coordinate  (C)
						(1,0)  coordinate (B)
						($(A)!  1/2 !(C)$) coordinate (O);
					\draw  (A)--(B)--(C)
						(A)--++(0,3.5)  coordinate (A')
						 (B)--++(0,3.5) coordinate (B')
						(C)--++(0,3.5) coordinate (C') 
						 (A')--(B')--(C')--cycle;
					\draw[dashed]  (B)--(O)--++(0,3.5) coordinate (D)
						 (A)--(C);
					 \draw[->] (D)--++(0,1) node[above]{$z$}; \draw[->] (C)--++(1,0) node[above]{$y$}; \draw[->] (B)--++(-0.5,-0.8) node[left]{$x$};
					\foreach \p/\r in {A/180,B/-60,C/-50,O/-50,A'/180,B'/-40,C'/0} 
					\fill  (\p) circle (1.2pt) node[shift={(\r:3mm)}]{$\p$};
				\end{tikzpicture}
		\end{center}
		Giả sử lăng trụ tam giác đều $ABC.A'B'C'$ có tất cả các cạnh bằng  $1$.\\
		Chọn  hệ trục $Oxyz$, với $O$ là trung điểm $AC$, $B \in Ox$, $C \in Oy$.\\
		Ta có $A\left(0;-\dfrac{1}{2};0\right)$,  $B\left(\dfrac{\sqrt{3}}{2};0;0\right)$, $C\left(0;\dfrac{1}{2};0\right)$, $B'\left(\dfrac{\sqrt{3}}{2};0;1\right)$, $C'\left(0;\dfrac{1}{2};1\right)$.\\
		$\overrightarrow{AB} = \left(\dfrac{\sqrt{3}}{2};\dfrac{1}{2};0\right)$, $\overrightarrow{BC'} = \left(-\dfrac{\sqrt{3}}{2};\dfrac{1}{2};1\right)$, $\overrightarrow{CC'} =  (0;0;1)$.\\
		 Mặt phẳng $\left( C'AB\right) $ có vectơ pháp tuyến $\overrightarrow{n}_1 = \left[\overrightarrow{AB}, \overrightarrow{BC'}\right] = \left(\dfrac{1}{2};-\dfrac{\sqrt{3}}{2};\dfrac{\sqrt{3}}{2}\right)$.\\
		Mặt phẳng $\left( BCC'B'\right) $ có vectơ pháp tuyến $\overrightarrow{n}_2 = \left[\overrightarrow{BC'}, \overrightarrow{CC'}\right]  = \left(\dfrac{1}{2};\dfrac{\sqrt{3}}{2};0\right)$.\\
		Ta có $\cos\alpha = \dfrac{\left|\overrightarrow{n}_1.\overrightarrow{n}_2\right|}{\left|\overrightarrow{n}_1\right|.\left|\overrightarrow{n}_2\right|} = \dfrac{\left|\dfrac{1}{2}\cdot\dfrac{1}{2} + \left(-\dfrac{\sqrt{3}}{2}\right)\cdot\dfrac{\sqrt{3}}{2} +  \dfrac{\sqrt{3}}{2}\cdot0\right|}{\sqrt{\dfrac{1}{4}+\dfrac{3}{4}+\dfrac{3}{4}}\cdot\sqrt{\dfrac{1}{4}+\dfrac{3}{4}+0}} = \dfrac{\sqrt{7}}{7}$.\\
		Ta lại có $1+\tan^2\alpha = \dfrac{1}{\cos^2\alpha} \Rightarrow \tan\alpha =  \sqrt{\dfrac{1}{\cos^2\alpha}-1} = \sqrt{\dfrac{1}{\left(\dfrac{\sqrt{7}}{7} \right)^2 }-1}=\sqrt{6}$.
	 }
\end{ex}

% nguồn: Phan1_Lop_12_Hinh_hoc_Chuong_5
\begin{ex}%[2H5V1-6]%[ma:fd46eeac]
	Trong không gian với hệ trục tọa độ $Oxyz$, cho hai mặt phẳng $(P)\colon \dfrac{x-2}{3}+\dfrac{y-1}{2}+\dfrac{z-4}{-6}=1$ và $(Q)\colon x+2y+3z+7=0$. Tính tang  góc tạo bởi hai mặt phẳng đã cho.
	 \choice
		 {$\dfrac{3}{\sqrt{19}}$}
		 {$\dfrac{3}{5\sqrt{19}}$}
		{$\dfrac{5}{3\sqrt{19}}$} 
		{\True $\dfrac{3\sqrt{19}}{5}$} 
	\loigiai{ 
		Mặt phẳng $(P)\colon \dfrac{x-2}{3}+\dfrac{y-1}{2}+\dfrac{z-4}{-6}=1\Leftrightarrow(P)\colon 2x+3y-z-9=0$.\\ 
		Mặt phẳng  $(P)$ có một vectơ pháp tuyến là $\overrightarrow{n}_{(P)}=(2; 3;-1)$.\\
		 $(Q)\colon x+2y+3z+7=0\Rightarrow \overrightarrow{n}_{(Q)}=(1; 2; 3)$.\\
		Gọi $\alpha$ là góc giữa hai mặt phẳng $(P)$ và  $(Q) \Rightarrow 0^\circ\leq \alpha \leq 90^\circ$.\\
		Ta có $\cos \alpha=\dfrac{\left|\overrightarrow{n}_{(P)} \cdot  \overrightarrow{n}_{(Q)}\right|}{\left|\overrightarrow{n}_{(P)}\right|\cdot\left|\overrightarrow{n}_{(Q)}\right|}=\dfrac{|2\cdot 1+3\cdot 2+(-1) \cdot 3|}{\sqrt{2^2+3^2+(-1)^2} \cdot \sqrt{1^2+2^2+3^2}}=\dfrac{5}{14}$.\\
		 $\tan^2\alpha=\dfrac{1}{\cos^2\alpha}-1=\dfrac{171}{25} \Rightarrow \tan \alpha=\dfrac{3\sqrt{19}}{5}$.
	} 
\end{ex}

% nguồn: Phan1_Lop_12_Hinh_hoc_Chuong_5
\begin{ex}%[2H5V1-6]%[ma:43f8bcc5]
	Cho hình chóp $S.ABC$  có đáy $ABC$ là tam giác đều cạnh $a$ và độ dài cạnh $SA=a$ và vuông góc với $(ABC)$. Góc giữa mặt phẳng $(SAB)$ và mặt phẳng $(SBC)$ là $\alpha$. Khi đó $\tan\alpha$ bằng
	\choice 
		 {$1$}
		 {$\sqrt{3}$}
		 {\True $\sqrt{6}$}
		{$\dfrac{1}{2\sqrt{2}}$} 
	\loigiai{ 
		\begin{center}
			\begin{tikzpicture}[line join = round,line cap=round,>=stealth,font=\footnotesize,scale=1]
					\path 
						(0,1.5)  coordinate (A)
						(1.5,0)  coordinate (B)
						(5,1.5) coordinate (C) 
						 (0,5) coordinate (S)
						($(B)!  1/2 !(C)$) coordinate (H);
					\draw  (A)--(B)--(C)
						 (S)--(A)
						 (S)--(B)
						(S)--(C); 
					\draw[dashed] (H)--(A)--(C); 
					\draw[->] (H)--++(0,5) node[right]{$z$}; \draw[->] (C)--(6,2) node[above]{$x$};  \draw[->] (A)--(-1,1.8) node[above]{$y$};
					\foreach \p/\r in  {A/-130,B/-90,C/-50,S/90,H/-50}
					\fill (\p)  circle (1.2pt) node[shift={(\r:3mm)}]{$\p$};
				\end{tikzpicture}
		\end{center}
		Gắn hệ tọa độ $Oxyz$ như hình  vẽ\\
		Ta có $H(0;0;0)$, $C\left(\dfrac{a}{2};0;0\right)$, $B\left(-\dfrac{a}{2};0;0\right)$, $A\left(0;\dfrac{a\sqrt{3}}{2};0\right)$, $S\left(0;\dfrac{a\sqrt{3}}{2};a\right)$.\\ 
		Ta có $\overrightarrow{AB}=\left(-\dfrac{a}{2}, -\dfrac{a\sqrt{3}}{2},0\right)$, $\overrightarrow{AS}=(0,  0, a)$, $\overrightarrow{BS}=\left(\dfrac{a}{2}, \dfrac{a\sqrt{3}}{2}, a\right)$, $\overrightarrow{BC}=(a, 0, 0)$.\\
		Suy  ra $\overrightarrow{n}_{(SAB)}=\left[\overrightarrow{AB},\overrightarrow{AS}\right]=\left(-\dfrac{a^2\sqrt{3}}{2};-\dfrac{a^2}{2};0\right)$, $\overrightarrow{n}_{(SBC)}=\left[\overrightarrow{BS},\overrightarrow{BC}\right]=\left(0;a^2;-\dfrac{a^2\sqrt{3}}{2}\right)$.\\
		 $\cos\alpha=\left|\dfrac{\overrightarrow{n}_{(SAB)}\cdot\overrightarrow{n}_{(SBC)}}{\left|\overrightarrow{n}_{(SAB)}|\cdot|\overrightarrow{n}_{(SBC)}\right|}\right|=\dfrac{\dfrac{a^4}{2}}{\sqrt{\dfrac{3a^4}{4}+\dfrac{a^4}{4}}\cdot\sqrt{a^4+\dfrac{3a^4}{4}}}=\dfrac{1}{\sqrt{7}}$.\\
		Vậy  $\tan\alpha = \sqrt{\dfrac{1}{\cos^2\alpha}-1} = \sqrt{6}$.
	 }
\end{ex}

% nguồn: Phan1_Lop_12_Hinh_hoc_Chuong_5
\begin{ex}%[2H5V1-6]%[ma:fc21ea8e]
	Cho hình chóp $S.ABCD$ có đáy $ABCD$ là hình vuông cạnh bằng $a$. Cạnh bên $SA=a$  vuông góc với đáy $(ABCD)$. Gọi $M$ là điểm nằm trên cạnh $SD$ sao cho $SM = 2MD$. Hãy tính cosin của góc tạo bởi hai mặt phẳng $(AMC)$ và $(SBC)$?
	 \choice
		{$\dfrac{1}{2}$} 
		{\True  $\dfrac{\sqrt{3}}{2}$}
		 {$\dfrac{2}{\sqrt{5}}$}
		{$1$} 
	 \loigiai{
		\begin{center}
			\begin{tikzpicture}[line join = round,line cap=round,>=stealth,font=\footnotesize,scale=1]
					\draw  (0,0) node[left]{$B (a; 0; 0)$}--(3.5,0) node[below]{$C (a; a; 0)$}--(5.5,1) node[above right]{$D (0; a; 0)$}
						(2,5) node[left]{$S (0;  0; a)$}--(0,0)
						(2,5)--(3.5,0) 
						 (2,5)--(5.5,1)
						 (3.5,0)--($(2,5)!2/3!(5.5,1)$) node[above right]{$M$};
					\draw[dashed]  (0,0)--(2,1) node[above left]{$A (0; 0; 0)$}--(5.5,1)
						(2,5)--(2,1) 
						 ($(2,5)!2/3!(5.5,1)$)--(2,1)--(3.5,0);
					\draw[->] (0,0)--++(-1.3,-0.7) node[left]{$x$};  \draw[->] (2,5)--++(0,1) node[right]{$y$}; \draw[->] (5.5,1)--++(1,0) node[right]{$y$};
				\end{tikzpicture}
		\end{center}
		Chọn hệ trục tọa độ $Oxyz$ sao cho điểm $A \equiv O$, các điểm $B$, $D$, $S$ lần lượt thuộc chiều dương các trục tọa độ $Ox$, $Oy$,  $Oz$. Suy ra tọa độ các điểm như trên hình vẽ.\\
		Do $M$ là điểm nằm trên cạnh $SD$ sao cho $SM = 2MD$ suy ra\\ 
		$\overrightarrow{SM} = \dfrac{2}{3}\overrightarrow{SD} \Rightarrow M\left(0;\dfrac{2a}{3};\dfrac{a}{3}\right)$ $\Rightarrow \heva{&\overrightarrow{AM}=\left(0;\dfrac{2a}{3};\dfrac{a}{3}\right)\\&\overrightarrow{AC}=(a;a;0).}\\ \Rightarrow \left[\overrightarrow{AM},\overrightarrow{AC}\right]=\left(-\dfrac{a^2}{3};\dfrac{a^2}{3};-\dfrac{2a^2}{3}\right)=-\dfrac{a^2}{3}(1;-1;2)$.\\  Ta có $\heva{&\overrightarrow{SB}=(a;0;-a)\\&\overrightarrow{SC}=(a;a;-a)}\Rightarrow \left[\overrightarrow{SB},\overrightarrow{SC}\right]=(a^2;0;a^2)=a^2(1;0;1)$.\\
		Chọn vectơ pháp tuyến  của mặt phẳng $(AMC)$ là $\overrightarrow{n}_1 = (1;-1;2)$.\\
		Chọn vectơ pháp tuyến của mặt phẳng  $(SBC)$ là $\overrightarrow{n}_2 = (1;0;1)$.\\
		Khi đó $$\cos\left((AMC),(SBC)\right) = \dfrac{\left|\overrightarrow{n}_1\cdot\overrightarrow{n}_2\right|}{\left|\overrightarrow{n}_1\right|\cdot\left|\overrightarrow{n}_2\right|} = \dfrac{|1\cdot1 + (-1)\cdot0 + 2\cdot1|}{\sqrt{1^2+(-1)^2+2^2}\cdot\sqrt{1^2+0^2+1^2}}= \dfrac{3}{\sqrt{6}\cdot\sqrt{2}} =  \dfrac{\sqrt{3}}{2}.$$
	 }
\end{ex}


%%%%%%%%%%%%%%%% ĐÚNG/SAI %%%%%%%%%%%%%%%%

% nguồn: Phan1_Lop_12_Hinh_hoc_Chuong_5
\begin{ex}%[2H5H1-6]%[ma:5c90c989]
	Trong không gian $Oxyz$, cho  hai mặt phẳng $(P)\colon 2x - y + z + 4 = 0$ và $(Q)\colon x + y + z - 2 = 0$.
	\choiceTF 
		{\True Điểm $A(-1;2;0)$ thuộc  mặt phẳng $(P)$}
		{\True Điểm  $B(2;1;-3)$ không thuộc mặt phẳng $(P)$}
		{Mặt phẳng $(P)$ có một  vectơ pháp tuyến là $\overrightarrow{n} = (2;1;1)$}
		{\True Cosin của góc tạo bởi hai mặt phẳng $(P)$ và $(Q)$ là $\dfrac{\sqrt{2}}{3}$} 
	\loigiai{\begin{itemchoice}
	\itemch \textbf{Đúng}. Thay tọa độ điểm $A(-1;2;0)$ vào phương trình $(P)$ ta được  $2\cdot(-1) - 2 + 0 + 4 = 0\Leftrightarrow0=0$ đúng.\\
	Vậy $A\in(P)$. 
	\itemch \textbf{Đúng}. Thay tọa độ điểm $B$ vào  phương trình $(P)$ ta được $2\cdot(2) - 1 + (-3) + 4 = 0\Leftrightarrow4=0$ sai.
	Vậy $B\not\in(P)$. 
	\itemch \textbf{Sai}. Mặt phẳng $(P)\colon 2x - y + z + 4  = 0$ có vectơ pháp tuyến $\overrightarrow{n}=(2;-1;1)$.
	\itemch \textbf{Đúng}. 
	Hai mặt phẳng $(P)$  và $(Q)$ có các vectơ pháp tuyến lần lượt là $\overrightarrow{n}_P = (2;-1;1)$ và $\overrightarrow{n}_Q = (1;1;1)$.\\
	Gọi $\phi$ là số  đo của góc giữa $(P)$ và $(Q)$, ta có\\
	$\cos \phi = \dfrac{\big|\overrightarrow{n}_P \cdot \overrightarrow{n}_Q\big|}{\big|\overrightarrow{n}_P\big| \cdot \big|\overrightarrow{n}_Q\big|} = \dfrac{\big|2\cdot1 + (-1)\cdot1 + 1\cdot1\big|}{\sqrt{2^2+(-1)^2+1^2}  \cdot \sqrt{1^2+1^2+1^2}} = \dfrac{2}{\sqrt{6}\sqrt{3}}=\dfrac{\sqrt{2}}{3}$.
\end{itemchoice}
 }
\end{ex}

% nguồn: Phan1_Lop_12_Hinh_hoc_Chuong_5
\begin{ex}%[2H5V1-6]%[ma:23c31a76]
	Trong không gian với hệ tọa độ $Oxyz$, cho hai mặt phẳng $\left(P_1\right) \colon  3x+4y+7=0$ và $\left(P_2\right) \colon 5x+12z+17=0$.
	\choiceTF 
		{Mặt phẳng $\left(P_1\right)$ có một vectơ pháp tuyến với tọa độ là $(3;4;7)$} 
		{Mặt phẳng $\left(P_{2}\right)$  có một vectơ pháp tuyến với tọa độ là $(5;12;17)$}
		{\True  Tích vô hướng của hai vectơ với tọa độ $(3;4;0)$ và $(5;0;12)$ bằng $15$}
		{\True Góc  giữa hai mặt phẳng $\left(P_1\right)$ và $\left(P_2\right)$ (làm tròn đến hàng đơn vị của độ) bằng $77^\circ$}
	\loigiai{ 
		\begin{itemchoice}
			\itemch Mặt phẳng $\left(P_1\right)$ có một vectơ pháp tuyến với tọa độ là $(3;4;0)$,  không phải $(3;4;7)$.
			\itemch Mặt phẳng $\left(P_2\right)$ có một vectơ pháp tuyến với toạ  độ là $(5;0;12)$, không phải $(5;12;17)$.
			\itemch  Tích vô hướng của hai vectơ với toạ độ $(3;4;0)$ và $(5;0;12)$ bằng $3\cdot 5+4 \cdot 0+0 \cdot 12=15$.
			\itemch  Côsin góc giữa hai mặt phẳng $\left(P_1\right)$ và $\left(P_2\right)$ bằng $\dfrac{15}{\sqrt{3^2+4^2+0^2} \cdot \sqrt{5^2+0^2+12^2}}=\dfrac{3}{13}$.\\
			Góc giữa hai mặt phẳng $\left(P_1\right)$ và $\left(P_2\right)$ (làm tròn đến hàng đơn vị của độ) bằng  $77^\circ$.
		\end{itemchoice}
	 }
\end{ex}


%%%%%%%%%%%%%%%% TRẢ LỜI NGẮN %%%%%%%%%%%%%%%%

% nguồn: Phan2_Lop_12_Hinh_hoc_Chuong_5
\begin{bt}%[2H5H1-6]%[ma:44749bf2]
	Trên một cánh đồng điện mặt trời, người ta đã thiết lập sẵn một hệ toạ độ $Oxyz$. Hai tấm pin năng lượng lần lượt nằm trong hai mặt phẳng $(P)\colon x+3z+1=0$ và $\break(Q)\colon 2x+z+7=0$. Tính góc (đơn  vị độ) giữa $(P)$ và $(Q)$.
	 \shortans{45}
	 \loigiai{
		Mặt phẳng $(P)$ có véc-tơ pháp tuyến là  $\overrightarrow{n}_P=(1;0;3)$.\\
		Mặt phẳng $(Q)$ có véc-tơ pháp tuyến là  $\overrightarrow{n}_Q=(2;0;1)$.\\
		Gọi $\alpha$ là góc giữa hai mặt phẳng  $(P)$ và $(Q)$ $(0\le \alpha\le 90)$, ta có
		 $\cos\alpha = \left|\cos\left(\overrightarrow{n}_P,\overrightarrow{n}_Q\right)\right| = \dfrac{\left|\overrightarrow{n}_P\cdot\overrightarrow{n}_Q\right|}{\left|\overrightarrow{n}_P\right|\cdot\left|\overrightarrow{n}_Q\right|}= \dfrac{\left|1\cdot 2+0\cdot 0+3\cdot 1=5\right|}{\sqrt{1^2+0^2+3^2}\cdot\sqrt{2^2+0^2+1^2}}=\dfrac{1}{\sqrt{2}}.$
		Vậy  $\alpha=45^\circ$.
	} 
\end{bt}

% nguồn: Phan2_Lop_12_Hinh_hoc_Chuong_5
\begin{bt}%[2H5H1-6]%[ma:b109c143]
	Cho hình chóp $S.ABCD$ có đáy $ABCD$ là hình vuông cạnh có cạnh bằng $a$ và  $SA \perp (ABCD)$. Khi $a = 2$, hãy tính độ dài cạnh $SA$ để góc tạo bởi $(SBC)$ và $(SDC)$ bằng $60^\circ$.
	\shortans{2} 
	\loigiai{ 
		Chọn  hệ trục như hình vẽ dưới đây.\\
		Gọi $M$  là trung điểm của $SC$.\\
		Đặt $SA=2m$,  $m>0$.
		\begin{center}
			\begin{tikzpicture}[scale=0.65,font=\footnotesize,>=stealth]
					 \path
						(0,0) coordinate (D) 
						(-3,-2) coordinate (A) 
						(5,0)  coordinate (C)
						($(A)+(C)-(D)$)  coordinate (B)
						($(D)!0.5!(B)$)  coordinate (O)
						($(A)+(0,6)$) coordinate (S) 
						($(O)+(0,3)$) coordinate (M); 
					\draw[->] (C)--(7,0.5)  node[below]{$y$};
					 \draw[->] (B)--(3,-3) node[below]{$x$};
					\draw[->] (M)--(1,4.5) node[left]{$z$}; 
					\draw  (B)--(C)--(S)--(B)--(A)--(S);
					 \draw[dashed] (S)--(D)--(C)--(A)--(D)--(B)
						(M)--(O); 
					\foreach \x/\g  in {C/90,D/180,A/-100,B/-90,S/90,M/30,O/-90}\draw (\x) circle (.04) +(\g:.6) node{\footnotesize$\x$};
				\end{tikzpicture}
		\end{center}
		Khi đó $B\left(  \dfrac{a}{\sqrt{2}}; 0; 0 \right),\
		D\left( -\dfrac{a}{\sqrt{2}};  0; 0 \right),\
		C\left( 0;  \dfrac{a}{\sqrt{2}}; 0 \right),\
		M(0;0;m),\  S(0;0;2m).$\\
		Ta có $ 
		 \vec{n}_{SBC} = \left( 1;1;\dfrac{a}{m\sqrt{2}} \right),\quad
		 \vec{n}_{SDC} = \left( -1;1;\dfrac{a}{m\sqrt{2}} \right).$\\
		Yêu cầu bài toán $\cos \theta = \dfrac{1}{2}$, với $\theta = \widehat{(SBC),(SDC)}$.\\ 
		Áp dụng công thức  tích vô hướng
		$$ 
		\cos\theta  = \dfrac{|\vec{n}_1 \cdot \vec{n}_2|}{|\vec{n}_1|\cdot\|\vec{n}_2\|}
		= \dfrac{| -1 +  1 + \left(\dfrac{a}{m\sqrt{2}}\right)^2 |}
		{1 + 1 + \left( \dfrac{a}{m\sqrt{2}}  \right)^2} = \dfrac{1}{2}
		\Rightarrow \left( \dfrac{a}{m\sqrt{2}} \right)^2 = 1 
		\Rightarrow \dfrac{a}{m\sqrt{2}}  = \pm 1.
		 $$
		Với $a = 2 \Rightarrow  m = \sqrt{2} \Rightarrow SA = 2m = 2$.
	} 
\end{bt}

% nguồn: Phan2_Lop_12_Hinh_hoc_Chuong_5
\begin{bt}%[2H5H1-6]%[ma:441798af]
	Cho hình chóp $S.ABCD$ có đáy $ABCD$ là hình vuông có độ dài đường chéo bằng $a\sqrt{2}$ và $SA$ vuông góc với mặt phẳng $(ABCD)$. Gọi $\alpha$ là góc giữa hai mặt phẳng $(SBD)$ và $(ABCD)$. Nếu $\tan\alpha = \sqrt{2}$ thì góc giữa hai mặt phẳng $(SAC)$  và $(SBC)$ bằng bao nhiêu độ?
	\shortans{60} 
	\loigiai{ 
		Gọi $I  = AC \cap BD$.\\
		Hình vuông $ABCD$ có độ dài đường chéo bằng $a\sqrt{2}$ suy ra hình  vuông có cạnh bằng $a$.\\
		Ta có $\left((SBD) \cap (ABCD)\right) =  BD$ nên $\left((SBD);(ABCD)\right) = (SI;AI) = \widehat{SIA}$.\\
		$\tan \widehat{SIA} = \dfrac{SA}{AI}  \Rightarrow SA = a$.\\
		 Chọn hệ tọa độ $Oxyz$ như hình vẽ.\\
		\begin{center}
			\begin{tikzpicture}[scale=0.6,font=\footnotesize,>=stealth]
					\path 
						(0,0) coordinate  (A)
						(-2,-2) coordinate (B) 
						 (5,0) coordinate (D)
						 ($(B)+(D)-(A)$) coordinate (C)
						($(A)+(0,3)$)  coordinate (S)
						(intersection of B--D  and A--C) coordinate (I);
					\draw[->]  (D)--(7,0) node[below]{$y$};
					\draw[->]  (B)--(-3,-3) node[below]{$x$};
					\draw[->] (S)--(0,4)  node[left]{$z$};
					 \draw (C)--(D)--(S)--(C)--(B)--(S);
					 \draw[dashed] (S)--(A)--(D)--(B)--(A)--(C);
					 \pic[draw,thin,angle radius=2mm] {right angle = B--A--D};
					\foreach \x/\g in  {A/180,B/-90,C/-100,D/-80,S/170,I/-90}\draw (\x) circle (.04) +(\g:.5) node{\footnotesize$\x$};
				\end{tikzpicture}
		\end{center}
		Ta có $A(0;0;0), B(a;0;0), C(a;0;a), S(0;0;a)$.\\ 
		Khi đó $\vec{SA} = (0;0;a),\ \vec{SC}  = (a;0;-a),\ \vec{SB} = (a;0;-a)$.\\
		Mặt phẳng $(SAC)$ có vectơ pháp  tuyến $\vec{n}_1 = (-1;0;1)$.\\
		Mặt phẳng $(SBC)$ có vectơ pháp  tuyến $\vec{n}_2 = (1;0;1)$.\\
		Suy ra $\cos((SAC);(SBC)) = \dfrac{|\vec{n}_1 \cdot \vec{n}_2|}{\left|\vec{n}_1\right|\cdot\left|\vec{n}_2\right|} = \dfrac{1}{\sqrt{2} \cdot \sqrt{2}} = \dfrac{1}{2} \Rightarrow \widehat{((SAC);(SBC))}  = 60^\circ$.
	 }
\end{bt}

% nguồn: Phan2_Lop_12_Hinh_hoc_Chuong_5
\begin{bt}%[2H5H1-6]%[ma:e8f365eb]
	Cho hình vuông $ABCD$ cạnh $a$. Trên hai tia $Bz$, $Dt$ vuông góc với mặt phẳng $(ABCD)$ và cùng chiều, lần lượt lấy hai điểm $M$, $N$ sao cho $BM = \dfrac{a}{4}$, $DN =  2a$. Tính góc $\varphi$ giữa hai mặt phẳng $(AMN)$ và $(CMN)$. (Tính theo đơn vị độ).
	\shortans{90} 
	\loigiai{ 
		Gắn hệ  trục tọa độ như hình vẽ
		\begin{center}
			\begin{tikzpicture}[scale=0.6,font=\footnotesize,>=stealth]
					\path 
						(0,0) coordinate (A) 
						 (-3,-2) coordinate (B)
						 (5,0) coordinate (D)
						 ($(B)+(D)-(A)$) coordinate (C)
						($(B)+(0,2.5)$) coordinate  (M)
						 ($(D)+(0,5)$) coordinate (N)
						($(B)!2!(A)$) coordinate (y) 
						 ($(B)!1.5!(C)$) coordinate (x)
						($(M)+(0,4)$) coordinate  (z);
					\draw[dashed][->]  (A)--(y) node[below]{$y$};
					 \draw[->] (C)--(x) node[below]{$x$};
					\draw[->] (M)--(z) node[left]{$z$}; 
					 \draw (C)--(N)--(D)--(C)--(B)--(M)--(N);
					\draw[dashed] (B)--(A)--(D) 
						(M)--(A)--(N); 
					\foreach \x/\g in {A/120,B/-90,C/-100,D/30,M/180,N/60}\draw (\x) circle  (.05) +(\g:.5) node{\footnotesize$\x$};
				\end{tikzpicture}
		\end{center}
		Khi đó $B(0;0;0)$,  $A(0;a;0)$, $C(a;0;0)$, $M\left(0;0;\dfrac{a}{4}\right)$, $N(a;a;2a)$.\\
		Ta có 
		$\vec{AM} = \left(0;-a;\dfrac{a}{4}\right)$, $\vec{AN} =  (0;0;2a)$.\\
		$\Rightarrow$ mặt phẳng $(AMN)$  có vectơ pháp tuyến là
		$\vec{n}_{1}  = [\vec{AM},\vec{AN}] = \left(-2a^2; -\dfrac{a^2}{4}; 0\right)$.\\
		 Tương tự ta có
		$\vec{CM} = \left(-a;0;\dfrac{a}{4}\right)$, $\vec{CN} =  (0;a;2a)$\\
		 $\Rightarrow$ mặt phẳng $(CMN)$ có vectơ pháp tuyến là
		$\vec{n}_{2} = [\vec{CM},\vec{CN}] = \left(-\dfrac{a^2}{4};  2a^2; -a^2\right)$.\\
		 Vậy
		 $$
		\cos\varphi  = \dfrac{|\vec{n}_1 \cdot \vec{n}_2|}{|\vec{n}_1| \cdot |\vec{n}_2|}
		= \dfrac{\left|-2a^2 \cdot \left(-\dfrac{a^2}{4}\right) + \left(-\dfrac{a^2}{4}\right) \cdot 2a^2\right|}{\sqrt{4a^4  + \dfrac{a^4}{16}} \cdot \sqrt{\dfrac{a^4}{16} + 4a^4 + a^4}} = 0
		\Rightarrow \varphi  = 90^\circ.
		$$ 
	} 
\end{bt}

% nguồn: Phan2_Lop_12_Hinh_hoc_Chuong_5
\begin{bt}%[2H5H1-6]%[ma:c6913e2f]
	Một máy bay  cất cánh tại một sân bay, sau khi bắt đầu cất cánh trong thời gian ngắn máy bay sẽ bay theo một đường thẳng và sân bay nơi máy bay cất cánh được coi là một mặt phẳng. Chọn hệ tọa độ $Oxyz$, đơn vị độ dài trên mỗi trục là $1$ km. Biết rằng, máy bay bắt đầu cất cánh tại điểm $O(0;0;0)$ và sau một thời gian ngắn máy bay bay đến điểm $A(2;5;1{,}2)$ và sân bay máy bay cất cánh nằm trên mặt phẳng $(Oxy)$. Góc tạo bởi đường bay của máy bay cất cánh và sân bay bằng bao nhiêu (\textit{làm tròn kết quả đến hàng đơn vị theo đơn vị độ})?
	\shortans{13} 
	\loigiai{ 
		Dựa vào đề bài ta có thể thấy góc tạo bởi đường bay của máy bay cất cánh và  sân bay chính là góc của $\overrightarrow{OA}$ và mặt phẳng $(Oxy)$.\\
		Ta có: $\overrightarrow{OA}=(2;5;1{,}2)$ và  vectơ pháp tuyến của mặt phẳng $(Oxy)$ là $\overrightarrow{n}=\overrightarrow{k}=(0;0;1)$.\\
		Gọi $\alpha$ là  góc tạo bởi $\overrightarrow{OA}$ và mặt phẳng $(Oxy)$, suy ra:
		\begin{align*}
			&  \sin\alpha=\dfrac{|\overrightarrow{OA}\cdot\overrightarrow{k}|}{|\overrightarrow{OA}||\overrightarrow{k}|}=\dfrac{|1{,}2\cdot1|}{\sqrt{2^2+5^2+(1{,}2)^2}\cdot\sqrt{1^2}}=\dfrac{6\sqrt{761}}{761}.
		\end{align*}
		Suy ra $\alpha\approx13^\circ.$ 
	 }
\end{bt}


%%%%%%%%%%%%%%%% TỰ LUẬN %%%%%%%%%%%%%%%%

% nguồn: Phan2_Lop_12_Hinh_hoc_Chuong_5
\begin{bt}%[2H5H1-6]%[ma:4690cc3a]
	Cho hai mặt phẳng $(P)\colon x-y-6=0$ và $(Q)$. Biết rằng điểm $H(2;-1;-2)$ là hình chiếu vuông góc của gốc toạ độ $O(0;0;0)$ xuống mặt phẳng $(Q)$. Tính góc  giữa mặt phẳng $(P)$ và mặt phẳng $(Q)$.
	 \loigiai{
		$(P)$ có vectơ pháp tuyến là $\overrightarrow{n}_1=(1; -1; 0)$, $(Q)$ có vectơ pháp tuyến là $ \overrightarrow  {OH}=(2; -1; -2)$.\\
		Gọi $\varphi$  là góc giữa hai mặt phẳng $(P)$ và $(Q)$ ta có $$\cos\varphi=\dfrac{|2\cdot 1+1\cdot 1-2\cdot 0|}{\sqrt{1^2+(-1)^2+0^2}\cdot\sqrt{2^2+(-1)^2+(-2)^2}}=\dfrac{\sqrt{2}}{2}\Rightarrow \varphi=45^{\circ}.$$
	 }
\end{bt}

% nguồn: Phan2_Lop_12_Hinh_hoc_Chuong_5
\begin{bt}%[2H5H1-6]%[ma:c68eeb40]
	Trong không gian $Oxyz$, cho hai mặt  phẳng $(P)\colon 2x-y-z-3=0$ và $(Q)\colon x-z-2=0$. Góc giữa hai mặt phẳng $(P)$ và $(Q)$ bằng bao nhiêu?
	\loigiai{ 
		Mặt  phẳng $(P)$ có một vectơ pháp tuyến $\overrightarrow{n}_P=(2;-1;-1)$.\\
		 Mặt phẳng $(Q)$ có một vectơ pháp tuyến $\overrightarrow{n}_Q=(1; 0;-1)$.\\
		Ta có $\cos[(P),(Q)]=\left|\cos \left(\overrightarrow{n}_P, \overrightarrow{n}_Q\right)\right|=\dfrac{\left|\overrightarrow{n}_P \cdot \overrightarrow{n}_Q\right|}{\left|\overrightarrow{n}_P\right|\cdot\left|\overrightarrow{n}_Q\right|}=\dfrac{|2+0+1|}{\sqrt{4+1+1} \cdot  \sqrt{1+1}}=\dfrac{\sqrt{3}}{2}$.\\
		Vậy $[(P),(Q)]=30^\circ$. 
	} 
\end{bt}

% nguồn: Phan2_Lop_12_Hinh_hoc_Chuong_5
\begin{bt}%[2H5H1-6]%[ma:b27a7b39]
	 Trong không gian $Oxyz$ cho hai đường thẳng $d_1\colon \dfrac{x}{-1}=\dfrac{y+1}{2}=\dfrac{z}{2}$, $d_2\colon \heva{&x=2t\\&y=1\\&z=1-t}$. Gọi $\phi$ là góc giữa hai đường thẳng $d_1$, $d_2$. Giá trị $\cos\phi$ có dạng $\dfrac{a\sqrt{c}}{b}$. Tính giá trị biểu thức $P=b-3a+c$?
	 \loigiai{
		Véc-tơ chỉ phương của hai  đường thẳng $d_1$, $d_2$ lần lượt là $\overrightarrow{u}_{d_1}=(-1;2;2)$, $\overrightarrow{u}_{d_2}=(2;0;-1)$.\\
		Khi đó  $\cos\phi=\dfrac{\left|\overrightarrow{u}_{d_1}\cdot \overrightarrow{u}_{d_2}\right|}{\left|\overrightarrow{u}_{d_1}\right|\left|\overrightarrow{u}_{d_2}\right|}=\dfrac{\left|-1\cdot 2+2\cdot 0+2\cdot (-1)\right|}{\sqrt{(-1)^2+2^2+2^2}\cdot\sqrt{2^2+0^2+(-1)^2}}=\dfrac{4}{3\sqrt{5}}=\dfrac{4\sqrt{5}}{15}$.\\
		Vậy $a=4$, $b=15$, $c=5\Rightarrow b-3a+c=15-3\cdot  4+5=8$.
	} 
\end{bt}

% nguồn: Phan2_Lop_12_Hinh_hoc_Chuong_5
\begin{bt}%[2H5H1-6]%[ma:b4edcbe5]
	Trong không gian $Oxyz$, cho ba điểm $M(1; 0; 0), N(0; 1;  0)$ và $P(0; 0; 1)$. Cô-sin của góc giữa hai mặt phẳng $(MNP)$ và mặt phẳng $(Oxy)$ bằng bao nhiêu?
	\loigiai{ 
		Mặt phẳng $(MNP)$ có một vectơ pháp tuyến là $\overrightarrow{n}=[\overrightarrow{MN},  \overrightarrow{MP}]=(1; 1; 1)$.\\
		Mặt phẳng $(Oxy)$ có  một vectơ pháp tuyến là $\overrightarrow{k}=(0; 0; 1)$.\\
		Gọi $\varphi$  là góc giữa hai mặt phẳng $(MNP)$ và $(Oxy)$.\\
		Ta có $\cos \varphi=|\cos(\overrightarrow{n}, \overrightarrow{k})|=\dfrac{|1\cdot 0+1\cdot 0+1\cdot  1|}{\sqrt{1^2+1^2+1^2}}=\dfrac{1}{\sqrt{3}}$.
	} 
\end{bt}

% nguồn: Phan2_Lop_12_Hinh_hoc_Chuong_5
\begin{bt}%[2H5V1-6]%[ma:cbce214d]
	Trên một cánh đồng điện  năng lượng mặt trời, người ta đã thiết lập sẵn một hệ tọa độ $Oxyz$. Hai tấm pin năng lượng lần lượt nằm trong hai mặt phẳng $(P)\colon 2x+2z+1=0$ và $\left(P'\right): x+z+7=0$.
	\begin{enumerate}
		\item Tính góc giữa $(P)$  và $\left(P'\right)$.
		\item  Tính góc hợp bởi $(P)$ và $\left(P'\right)$ với mặt đất $(Q)$ có phương trình $z=0$.
	\end{enumerate}
	 \loigiai{
		\begin{enumerate}
			\item Mặt phẳng $(P)$ có vectơ pháp tuyến là $\overrightarrow{n}=(2; 0;  2)$ và mặt phẳng $\left(P'\right)$ có vectơ pháp tuyến là $\overrightarrow{n}'=(1; 0; 1)$ nên $\cos \left((P),\left(P'\right)\right)=\dfrac{|2\cdot1+0\cdot0+2\cdot1|}{\sqrt{2^2+0^2+2^2} \cdot \sqrt{1^2+0^2+1^2}}=\dfrac{4}{4}=1$.\\
			Suy  ra $\left((P),\left(P'\right)\right)=0^\circ$.
			\item Mặt phẳng $(Q)$ có vectơ pháp tuyến là  $\overrightarrow{n}_Q=(0; 0; 1)$
		\end{enumerate}
		 \allowdisplaybreaks
		\begin{eqnarray*}
			&&\cos((P),(Q))=\dfrac{|2 \cdot 0+0 \cdot 0+2 \cdot 1|}{\sqrt{2^2+0^2+2^2} \cdot  \sqrt{1^2}}=\dfrac{2}{2\sqrt{2}}=\dfrac{1}{\sqrt{2}} \Rightarrow((P),(Q))=45^\circ. \\
			&&\cos \left(\left(P'\right),(Q)\right)=\dfrac{|1 \cdot 0+0 \cdot 0+1 \cdot 1|}{\sqrt{1^2+0^2+1^2} \cdot \sqrt{1^2}}=\dfrac{1}{\sqrt{2}} \Rightarrow\left(\left(P'\right),(Q)\right)=45^\circ. 
		\end{eqnarray*}
	 }
\end{bt}
